\begin{textsample}{POS Dim 2 – vem\_ed\_17 – Score 3.00 – vem\_ed\_17/t078.txt}  \label{ex:f2_pos_125}
A VISÃO DOS DOCENTES: PESQUISA 2º SEMESTRE 2022 Todos os 44 professores que \textbf{responderam} à pesquisa de percepção autorizaram \textbf{divulgar} as respostas para fins de publicações gerenciais e acadêmicas e todos recomendaram um PCI para colegas, como ocorreu nas pesquisas anteriores.
Quase dois terços (64 \%) são veteranos, realizando o PCI pela terceira vez ou mais: isso mostra o amadurecimento das colaborações internacionais nas Fatecs.
Na edição anterior da pesquisa (primeiro semestre de 2022), 49 \% dos professores respondentes realizavam o PCI pela primeira vez.
Aquele foi um momento de captação, por parte da equipe dos PCIs / Cesu, de novos docentes interessados, que mantiveram suas colaborações e aos poucos ganharam experiência internacional.
Essa consolidação se \textbf{mantém} para o futuro próximo: dentre os respondentes do segundo semestre de 2022, 66 \% professores pretendem continuar com mesmo projeto e parceiro, 23 \% continuam com o mesmo parceiro, apenas mudando o projeto, 7 \% vão buscar novo parceiro e apenas 4 \% não desejam realizar outro PCI.
Dos 44 pesquisados, 77,3 \% reconhecem que a competência em língua estrangeira melhorou.
É quase unânime (97,8 \%) que os PCIs possibilitam desenvolver novas atividades acadêmicas a partir da experiência do parceiro internacional.
Além disso, 93 \% concordam que os PCIs abrem portas para publicar pesquisas em sua área de atuação e 86,3 \% afirmam que podem elaborar artigos acadêmicos com o parceiro.
Esses dados \textbf{mantêm} a tendência positiva dos anos anteriores.

% matched lemmas: divulgar, manter, responder
\end{textsample}
